% 表紙・凡例・目次（記入例・模範解答版）
\thispagestyle{cover}
\vspace*{12mm}
{\sffamily\color{pone}\large 名古屋市立大学 経済学部　2026年度後期　水曜3限}\par
\vspace{3mm}
{\sffamily\bfseries\fontsize{30pt}{38pt}\selectfont 基礎演習Ⅱ}\par
\vspace{2mm}
{\sffamily\fontsize{20pt}{26pt}\selectfont 各回ワークシート　\textcolor{ans}{記入例・模範解答}}\par
\vspace{2mm}
{\sffamily\small\color{inktwo}教員用。学生への配布は、授業での扱いを決めてから。}\par
\vspace{3mm}
{\color{inktwo}\rule{\linewidth}{0.6pt}}\par
\vspace{6mm}

\begin{tcolorbox}[enhanced,colback=papertwo,colframe=papertwo,boxrule=0pt,arc=1.5mm,
    left=4mm,right=4mm,top=3mm,bottom=3mm,fontupper=\small]
{\sffamily\bfseries この冊子の見方}\par\vspace{1mm}
\begin{itemize}[label=・]
  \item 各回のワークシートに、\ans{赤字}で書き込みを入れています。
  \item 答えが一つに定まる設問（授業の仕組み、良い問いの条件、評価の観点など）は\textbf{模範解答}です。
  \item 学生ごとに答えが変わる設問（テーマ、問い、ペルソナ、コンセプトなど）は\textbf{記入例}です。下の架空の学生とグループが書いたものとして、半年分を一貫させています。正解ではなく、記述の粒度と具体性の目安として使ってください。
  \item 各回の末尾の「\textcolor{ans}{指導のポイント・解説}」に、評価の観点とつまずきやすい点をまとめています。
\end{itemize}
\end{tcolorbox}

\vspace{4mm}
\begin{tcolorbox}[enhanced,colback=white,colframe=rulegray,boxrule=0.5pt,arc=1.5mm,
    left=4mm,right=4mm,top=3mm,bottom=3mm,fontupper=\small]
{\sffamily\bfseries 記入例の前提（架空）}\par\vspace{1mm}
\begin{itemize}[label=・]
  \item \textbf{学生}：経済学部2年。実家の近くのカーディーラーが閑散としているのを見て、若い世代と車の関係に関心を持っている。
  \item \textbf{前半の問い}：「車の購入を検討した名古屋都市圏の20代は、来店前にオンラインで何を調べ、何を目的にディーラーへ来店しているのか」（第2回の仮の問いを、第5回で修正して確定）。
  \item \textbf{後半の課題}：郊外型カーディーラー（架空のＡ社）の店舗リニューアル提案。実際の課題は第6回に企業から提示されるため、記入例の課題設定・数値・社名はすべて架空です。
  \item \textbf{文献}：記入例に挙げた統計・白書は実在する継続調査です。版や年次、URLは配布前に最新のものを確認してください。第4回で精読する論文は、記入例のための\textbf{架空の論文}です。
\end{itemize}
\end{tcolorbox}

\vspace{6mm}
{\sffamily\bfseries 目次}\par\vspace{1mm}
{\small\makeatletter\@starttoc{toc}\makeatother}
